\documentclass{article}
\usepackage[margin=2cm]{geometry}
\usepackage{amsmath}
\usepackage{graphicx}
\usepackage{booktabs}
\usepackage[utf8]{inputenc}
\usepackage{ragged2e}
\setlength{\emergencystretch}{3em}
\begin{document}
\sloppy

\noindent\hfill\includegraphics[height=1.8em]{assets/figure_header_logo.png}\par\noindent\hrule\vspace{4pt}

The non-zero components of the Ricci tensor are


\begin{equation}
\begin{array} { c c c } { R _ { u u } = - 5 \left ( f ^ { \prime \prime } + ( f ^ { \prime } ) ^ { 2 } \right ) ~ R _ { m n } = - g _ { m n } \left ( f ^ { \prime \prime } + 5 ( f ^ { \prime } ) ^ { 2 } \right ) ~ }\end{array}
\end{equation}


while the Ricci scalar is given by


\begin{equation}
R = - 1 0 f ^ { \prime \prime } - 3 0 ( f ^ { \prime } ) ^ { 2 } \qquad ( 2 )
\end{equation}


Furthermore, we have that G = e5f g, where g is the determinant of the unit S metric.
Upon integration by parts, part of the Einstein-Hilbert term cancels with the Gibbons-
Hawking term to give the following simple expression


\begin{equation}
S = \int d u \int d ^ { 5 } x \sqrt { g } ~ e ^ { 5 f } \left [ - 5 \left ( ( f ^ { \prime } ) ^ { 2 } + 4 W ^ { 2 } \right ) + 2 \mathcal { L } _ { \mathrm { k i n } } \right ] \qquad \qquad \qquad \qquad ( 3 ) \qquad
\end{equation}


The restriction to the flat case was not strictly necessary so far, but it will be crucial in the
next step. The gradient flow equations, together with the chain-rule, allows us to rewrite


\begin{equation}
\mathcal { L } _ \mathsf { k i n } = - 2 ~ W ^ { \prime } \qquad ( 4 )
\end{equation}


Plugging this into and using the BPS equation of the warp factor, we find


\begin{equation}
S = - 4 \int d ^ { 6 } \! x \sqrt { g } ~ e ^ { 5 f } W \Big | _ { 6 } ^ { \Lambda } \qquad ( 5 )
\end{equation}


where A is the UV cutoff. Only the A part of the action contributes, since e5fW|o vanishes.
due to the close-off of the geometry. Removing the UV cutoff A  oo is equivalent to
removing the cutoff e on our asymptotic coordinate z, i.e. e -> 0. From the UV asymptotics
we find that in this limit the factor e5f diverges like


\begin{equation}
e ^ { 5 f } \sim \frac { 1 } { \varepsilon ^ { 5 } } \qquad ( 6 )
\end{equation}


tial are relevant for obtaining counterterms. All the higher-order terms vanish as the cutoff
is removed. We may thus legitimately insert the approximate superpotential into to get the
counterterms,


\begin{equation}
S _ { \mathrm { c l } } ^ { ( W ) } = 4 \int d ^ { 6 } x \sqrt { \gamma } \left [ \frac { 1 } { 2 } + \frac { 3 } { 4 } \sigma ^ { 2 } + \frac { 1 } { 1 6 } ( \phi ^ { 0 } ) ^ { 2 } - \frac { 3 } { 1 6 } ( \phi ^ { 3 } ) ^ { 2 } + \frac { 1 } { 1 9 2 } ( \phi ^ { 6 } ) ^ { 4 } - \frac { 3 } { 1 6 } ( \phi ^ { 0 } ) ^ { 2 } \sigma \right ] ~ ( 7 )
\end{equation}


where  is the induced metric on the z = e boundary. All fields are evaluated at z = e. This
gives all finite and infinite counterterms for the flat domain wall solutions..

\end{document}
